% 
%
% (c) 2025 Autor, ETH Zürich
%
% !TEX root = linalg_I-II_summary.tex
% !TEX encoding = UTF-8
%

\section{Spectral Theorems}

\begin{definition}{Orthogonal Diagonalizability}{orth_diag}
	Let $V$ be an inner product space and $T:V \to V$ and endomorphism. We say that $T$ is \textbf{orthogonally diagonalizable} if there exists an ortho\textcolor{red}{normal} basis for $V$ consisting of eigenvectors of $T$.
\end{definition}

\begin{lemma}{Consequence of Orthogonal Diagonalizability}{consq_orth_diag}
	If $T$ is orthogonally diagonalizable, then 
	\[
		T\circ T^* = T^* \circ T.
	\]
\end{lemma}

\begin{definition}{\textcolor{red}{Normal Endomorphism/Matrix}}{normal_endo}
	Let $V$ be an inner product space. An endomorphism $T:V \to V$ is called \textbf{normal} if 
	\[
		T \circ T^* = T^* \circ T.
	\]
	A matrix $A \in M_{n \times n}(K)$ ($K = \mathbb{R}, \mathbb{C}$) is called \textbf{normal} if 
	\[
		A A^* = A^* A.
	\]
	This is equivalent to saying that $T_A:K^n \to K^n$ is normal, if we endow $K^n$ with its standard inner product. 
\end{definition}

We've seen in Lemma \ref{lem*consq_orth_diag} that if $T:V \to V$ is orthogonally diagonalizable, then $T$ is normal. Is the converse statement also true?

Over $K = \mathbb{R}$ the answer is \textbf{no!}, e.g., 
\[
	A = \begin{pmatrix}
		0 & 1\\
		-1 & 0
	\end{pmatrix}.
\]
$T_A:\mathbb{R}^2 \to \mathbb{R}^2$ is a clockwise rotation by 90°, so there exist no eigenvalues in $\mathbb{R}$. But $AA^* = A^*A = I$.

\begin{theorem}{Spectral Theorem over $\mathbb{C}$}{spec_theo_c}
	Let $V$ be a finite dimensional inner product space over $\mathbb{C}$ and $T:V \to V$ a linear map. Then $T$ is orthogonally diagonalizable if and only if $T$ is normal.
\end{theorem}

\begin{lemma}{\textcolor{red}{Properties of a Normal Endomorphism}}{prop_normal_endo}
	Let $V$ be a finite dimensional inner product space over $K = \mathbb{R}$ or $\mathbb{C}$. Let $T:V \to V$ be a normal endomorphism. Then:
	\begin{enumerate}
		\item $\forall\,  v \in V : \; \Vert Tv \Vert = \Vert T^*v \Vert$.
		\item $\forall \, \lambda \in K$ the endomorphism $T - \lambda \cdot id_V$ is also normal.
		\item If $v$ is an eigenvector of $T$ with eigenvalue $\lambda$, then $v$ is also an eigenvector of $T^*$ with eigenvalue $\overline{\lambda}$.
		\item Let $v_1, v_2$ be eigenvectors of $T$ with eigenvalues $\lambda_1, \lambda_2$ respectively. If $\lambda_1 \neq \lambda_2$, then $\langle v_1, v_2\rangle = 0$.
	\end{enumerate}
\end{lemma}

\begin{proof}
	\textbf{1)} Let $v \in V$. Then by the properties of the adjoint map $T^*$ and since $T$ is normal, we have that
	\begin{align*}
		\Vert T v \Vert ^2 &= \langle Tv, Tv\rangle = \langle v, T^*\circ Tv\rangle = \langle v, T \circ T^* v\rangle \\
		&= \overline{\langle T \circ T^*v, v\rangle} = \overline{\langle T^*v, T^*v \rangle} = \langle T^*v, T^*v \rangle = \Vert T^* v\Vert ^2.
	\end{align*}
	Since the norm is non-negative, we conclude that
	\[
		\Vert Tv \Vert = \Vert T^*v\Vert.
	\]
	
	\textbf{2)} Let $\lambda \in K$. By properties of adjoint maps and the normality of $T$ we have that
\end{proof}